\section{Fluid mechanics and the load required by a target}\label{sec:mechanics}

\subsection{A continuum reference problem}
For each phase $\ell$, define the material derivative
$D_t=\partial_t+\bm v_\ell\cdot\nabla$, the deformation-rate tensor
$\bm D(\bm v)=(\nabla\bm v+\nabla\bm v^T)/2$, and its deviatoric part
$\bm D_0=\bm D-(\nabla\cdot\bm v)\Id/3$. Let $\bm f_\ell$ be an additional
prescribed body-force density, $\bm q_\ell$ the conductive heat flux, and
$r_\ell$ an externally supplied volumetric heating rate. Their units are
$\mathrm{N\,m^{-3}}$, $\mathrm{W\,m^{-2}}$, and $\mathrm{W\,m^{-3}}$,
respectively. A compressible Newtonian reference model is
\begin{align}
 \partial_t\rho_\ell+\nabla\cdot(\rho_\ell\bm v_\ell)&=0,
 \label{eq:mass}\\
 \rho_\ell D_t\bm v_\ell
 &=\nabla\cdot\bm T_\ell-\rho_\ell\nabla\Phi+\bm f_\ell,
 \label{eq:momentum}\\
 \rho_\ell D_te_\ell
 &=-p_\ell\nabla\cdot\bm v_\ell
   +\bm\tau_\ell:\nabla\bm v_\ell-\nabla\cdot\bm q_\ell+r_\ell,
 \label{eq:internal-energy}\\
 \bm T_\ell&=-p_\ell\Id+\bm\tau_\ell,\qquad
 \bm\tau_\ell=2\mu_\ell\bm D_0(\bm v_\ell)
                +\zeta_\ell(\nabla\cdot\bm v_\ell)\Id,
 \qquad \bm q_\ell=-k_\ell\nabla\Theta_\ell.
 \label{eq:constitutive}
\end{align}
The colon denotes tensor contraction. Equations of state
$p_\ell=p_\ell(\rho_\ell,\Theta_\ell)$ and
$e_\ell=e_\ell(\rho_\ell,\Theta_\ell)$ close the thermodynamics. Transport
coefficients may depend on the local state. Resolving the acoustic waveform in
these equations already accounts for its viscous heating and momentum transfer;
one must not then add a second radiation-force term for the same wave.
Thermoviscous perturbation models derive from these balances, with a consistent
separation of orders \cite{jorgensen2021}.

Let $V_\Gamma$ denote interface normal speed. Without mass transfer,
$V_\Gamma=\bm v^-\cdot\nn=\bm v^+\cdot\nn$. The clean-interface stress balance
is
\begin{equation}
 \jump{\bm T}\nn=\sigma\kappa\nn-\Gs\sigma.
 \label{eq:instantaneous-jump}
\end{equation}
For two viscous fluids the usual clean, no-interfacial-slip model also imposes
velocity continuity. Other interfacial constitutive laws must be specified if
surface viscosity, surfactants, or slip are admitted. With constant surface
tension, a static spherical inclusion obeys
$p^- -p^+=2\sigma/R_s$, where $R_s$ is its radius; this checks the sign in
\cref{eq:instantaneous-jump}.

A thermally passive interface with no thermal resistance or surface heat
storage has continuous temperature and normal conductive heat flux. If
$\sigma$ varies with temperature or composition, the associated surface
thermodynamics and species transport must accompany the Marangoni term in
\cref{eq:instantaneous-jump}. Keeping that force while claiming an exact
constant-$\sigma$ energy identity would be inconsistent.

The remaining boundaries require mechanical and thermal conditions: for
example prescribed motion, an elastic solid coupled by traction and velocity,
or a radiating exterior. These alternatives define different wave and flow
operators. Neither a rigid reflecting chamber nor a pressure-release free
surface is a universal boundary condition.

\subsection{Wetting is part of the problem data}
If a supporting solid is present, let $\theta$ denote contact angle measured
through the lens fluid, and let $\gamma_{s-}$ and $\gamma_{s+}$ be the solid's
interfacial free energies with the two phases, in $\mathrm{N\,m^{-1}}$. An ideal
equilibrium contact angle $\theta_e$ satisfies Young's relation
\begin{equation}
 \sigma\cos\theta_e=\gamma_{s+}-\gamma_{s-}.
 \label{eq:young-contact}
\end{equation}
A fixed contact line is a positional constraint; it is not equivalent to
prescribing this angle. For a mobile line, its law must be supplied. As one
explicit dissipative example, define $V_{\mathcal C}$ positive when the lens
fluid advances over the wall, and a line-friction coefficient
$\zeta_{\mathcal C}\geq0$ in $\mathrm{Pa\,s}$. Then
\begin{equation}
 \zeta_{\mathcal C}V_{\mathcal C}
   =\sigma(\cos\theta_e-\cos\theta)
 \label{eq:contact-friction}
\end{equation}
is a possible closure, not a universal wetting law. Hysteresis can instead
allow a stationary line over a measured interval of angles. The classical
moving-line/no-slip singularity means that imposing motion also requires a
compatible microscopic regularization, such as slip or a diffuse-interface
model \cite{huh1971,qian2006}. Neither a guessed contact angle nor a visually
retained meniscus replaces these data.

\subsection{Variational derivation of the required static traction}
Now restrict attention to constant-density fluids, constant $\sigma$, a fixed
support, and a quiescent state with no body force beyond $-\rho_\ell\nabla\Phi$.
The varied lens region has finite volume and deformable interfacial area. Set
$\Delta\rho=\rho_- -\rho_+$. Let $A[\Gamma]$ denote deformable interfacial area,
$V[\Gamma]$ the lens volume, and $E_{\rm wall}$ the solid wetting energy. Up to
constants independent of the shape, the mechanical potential energy is
\begin{equation}
 E[\Gamma]=\sigma A[\Gamma]
     +\Delta\rho\int_{\Omega^-}\Phi\dd V+E_{\rm wall}.
 \label{eq:capillary-energy}
\end{equation}
This is the standard capillary variational structure \cite{finn1986}. For a
virtual displacement $\eta\nn$ compatible with the support, area and volume
variations give
\begin{equation}
 \delta V=\int_\Gamma\eta\dd A,\qquad
 \delta E=\int_\Gamma(\sigma\kappa+\Delta\rho\Phi)\eta\dd A
                +\text{contact-line terms}.
 \label{eq:first-variation}
\end{equation}
The line terms either vanish for a pinned boundary or supply the appropriate
wetting condition. A normal acoustic traction $\Pi$, positive outward, does
virtual work $\int_\Gamma\Pi\eta\dd A$. Enforcing the prescribed volume
$V_0$ with the multiplier $\lambda_V$ therefore gives
\begin{equation}
 \boxed{\Pi_{\rm req}[\Gamma]
       =\sigma\kappa+\Delta\rho\Phi-\lambda_V,
       \qquad V[\Gamma]=V_0.}
 \label{eq:required-traction}
\end{equation}
Here $V_0$ is a supplied volume, not a fitted correction to the target. Every
term in the traction equation is a pressure. It determines the load required
by any admissible smooth target, whether Cartesian or otherwise.

For a graph $z=h(x,y)$ with the lens fluid below it and an upward normal, let
$\nabla_{xy}$ act in the horizontal plane. In a uniform downward gravitational
field with magnitude $g_{\rm grav}>0$, $\Phi=g_{\rm grav}z$, and
\begin{equation}
 \kappa=-\nabla_{xy}\cdot
       \frac{\nabla_{xy}h}{\sqrt{1+|\nabla_{xy}h|^2}},\qquad
 \Pi_{\rm req}=-\sigma\nabla_{xy}\cdot
       \frac{\nabla_{xy}h}{\sqrt{1+|\nabla_{xy}h|^2}}
       +\Delta\rho g_{\rm grav}h-\lambda_V.
 \label{eq:graph-traction}
\end{equation}
No small-slope approximation is made. Overhangs require a surface
representation, using \cref{eq:required-traction} directly.

Define the zero-mean projection of a surface scalar $f$ by
$\mathcal P_0f=f-A[\Gamma]^{-1}\int_\Gamma f\dd A$. For closed-volume shape
variations, only the nonconstant part of the normal load determines shape:
\begin{equation}
 \mathcal P_0\Pi=\mathcal P_0(\sigma\kappa+\Delta\rho\Phi).
 \label{eq:pressure-gauge}
\end{equation}
The omitted constant is balanced by the pressure multiplier. This gauge does
not remove restrictions on absolute pressure, cavitation, or actuator loading.
For several independently conserved volumes there is one corresponding
constraint per volume. A reservoir changes this formulation according to its
specified pressure and mass exchange.

The quiescent reduction additionally requires compatible tangential traction
and no uncompensated mean bulk forcing. If acoustic streaming produces a
stationary flow $\bm U_\star\ne0$, its pressure and viscous stresses must be
solved together with $\bm U_\star\cdot\nn=0$. Matching
\cref{eq:required-traction} alone does not establish that steady state.

\subsection{What a finite pulse can leave behind}
\begin{proposition}[Necessary condition for unforced retention]
Suppose actuation ends, all acoustic, flow and thermal transients decay, the
liquid remains liquid with the same constant densities and surface tension,
and the support, wetting law, volume and conservative body force are unchanged.
Any smooth quiescent limiting interface must satisfy
\begin{equation}
 \sigma\kappa+\Delta\rho\Phi=\lambda_V
 \label{eq:unforced-retention}
\end{equation}
with its volume and boundary conditions. A target violating this condition
cannot be a permanently retained quiescent result of that pulse.
\end{proposition}
\begin{proof}
At a quiescent limit, hydrostatic balance gives
$p^h_- -p^h_+=\lambda_V-\Delta\rho\Phi$. With no residual wave traction,
\cref{eq:instantaneous-jump} requires $p^h_- -p^h_+=\sigma\kappa$.
Equivalently, the constrained first variation of \cref{eq:capillary-energy}
vanishes. Both give \cref{eq:unforced-retention}.
\end{proof}
This is a necessary condition for a limit, not a proof that every trajectory
converges or that equilibrium is unique. Persistent rotation or another
non-quiescent relative equilibrium requires retaining that flow and its
conserved quantities; it is outside the stated quiescent-limit hypothesis.
Contact-angle hysteresis can leave a
different admissible pinned boundary. Nevertheless, it cannot support an
arbitrary non-capillary interior shape in a fully relaxed liquid. Continuous
actuation, a modified support, or curing can alter the retention question.
Fluidic fabrication by shaping and solidification is an experimentally
demonstrated but distinct route \cite{elgarisi2021}; it is not evidence that
an arbitrary shape persists after an acoustic pulse in an unchanged liquid.
